% ---------------------------------------------------------------------------
%  II. Related Work
%  Floats:  tables/tab-positioning.tex   (Table I)
%  Rule:    every subsection ends by saying what that line of work assumes
%           and why our starting point is different. Add a citation here and
%           it must also gain a row in Table I, or it is decoration.
% ---------------------------------------------------------------------------

\section{Related Work}
\label{sec:related}

\subsection{Re-engineering commercial robot platforms}
The closest precedent is Kim \emph{et~al.}~\cite{kim2005reengineering}, who re-engineered the
software architecture of a commercial home service robot. Their study follows much the same arc
as ours, but one difference decides everything: they had the manufacturer's source code and its
engineering team, while we start from a sealed product. M\"uhe \emph{et~al.}~\cite{muhe2010kuka}
recovered a proprietary industrial robot language, again working from artifacts the vendor had
published. Both begin with more material than we do, and the methodology below exists mainly to
make up for that shortfall.

\subsection{Vendor lock-in and interoperability}
The trade-off between proprietary and open platform strategies is long
established~\cite{west2003open}, and service robotics reproduces it one product at a time.
Glauser \emph{et~al.}~\cite{glauser2023interop} report healthcare deployments held back by
outdated and restrictive vendor SDKs, and propose an interoperable programming interface as the
remedy. Middleware bridges such as ROS--YARP~\cite{aragao2016yarp} address the adjacent problem
of connecting two interfaces that are both \emph{known}. In every case the interface itself is
taken for granted. The question we take up comes earlier: how does a laboratory obtain one at
all, on a robot it has already bought?

\subsection{Black-box access to \ros{} systems}
\rosbridge{} was designed to let non-\ros{} clients reach a \ros{} graph over
WebSocket~\cite{crick2012rosbridge,rosbridge_suite}, normally deployed by the integrator who
controls that graph. DeMarinis \emph{et~al.}~\cite{demarinis2019scanning} inverted this
perspective, scanning the Internet for exposed \ros{} systems and showing that sensors and
actuators can be reached without any knowledge of the node graph; Dieber
\emph{et~al.}~\cite{dieber2017security} treat the same openness as an application-level
security problem.
We occupy a similar position with the sign reversed: we too arrive at a vendor-deployed bridge
as an unknown client, but constructively, and on hardware we own.

\subsection{Protocol recovery from companion software}
Extracting protocol specifications from observed traffic is a mature security
technique~\cite{comparetti2009prospex}, and Android analysis has focused on information flow
within an application~\cite{enck2014taintdroid} or between applications~\cite{chin2011android}.
The practice of reading a device's companion app to recover its wire protocol is established in
the IoT community~\cite{giese2018xiaomi}. Hypothesis H1 imports that practice into robotics,
where the companion application controls a second, physically separate computer.

\subsection{Edge robotics and on-device interaction}
Cloud robotics assumes a reliable path to remote compute~\cite{kehoe2015cloud}. We had the
opposite constraint. The robot's own network guarantees no Internet route, so speech and
recognition had to run on the head, which pushed us towards closed-vocabulary
recognition~\cite{warden2018speechcommands,vosk} and on-device face
embeddings~\cite{schroff2015facenet}. Language-model agents~\cite{ahn2022saycan,driess2023paleme}
assume an open action API is already available. On a closed robot that API has to be recovered
first, which is the work reported here.

% Table I — positioning, by what each route presupposes.
% Used by sections/02-related-work.
%
% NOTE ON FRAMING. This table merges two that used to be separate: a
% related-work grid (Work / Emphasis / Distinction) and a route comparison.
% They answered the same question — how is this different from what exists —
% for papers and for approaches respectively, and the page budget does not
% carry both. Rows are ROUTES to a control interface, each citing its
% representative work, because that is the axis on which they actually differ.
%
% A latency column was deliberately not included: rosbridge is the transport
% this work runs on, not a competitor, so comparing throughput would compare a
% thing with itself. What separates these routes is what they require the
% vendor to have provided first, which is the last column.
\begin{table}[!t]
\caption{Routes to a control interface on a closed service robot, and what each presupposes}
\label{tab:positioning}
\centering
\scriptsize
\setlength{\tabcolsep}{3pt}
\begin{tabular}{@{}>{\raggedright\arraybackslash}p{1.75cm}ccc>{\raggedright\arraybackslash}p{2.3cm}@{}}
\toprule
\textbf{Route} & \textbf{Offl.} & \textbf{Insp.} & \textbf{N.-d.} & \textbf{Presupposes} \\
\midrule
Vendor SDK~\cite{glauser2023interop} & varies & no & yes &
A published SDK, kept current by the vendor \\[2pt]
Vendor cloud API & no & no & yes &
A subscription and a network path off site \\[2pt]
Bridge integration~\cite{crick2012rosbridge,aragao2016yarp} & yes & yes & yes &
Control of the \ros{} graph, and known message definitions \\[2pt]
Re-architecting with vendor support~\cite{kim2005reengineering} & yes & yes & yes &
Source code and access to the vendor's engineers \\[2pt]
Companion-app recovery~\cite{giese2018xiaomi,muhe2010kuka} & --- & yes & yes &
An artifact that already speaks the protocol \\[2pt]
Firmware replacement & yes & yes & \textbf{no} &
Root access; voids warranty, unrecoverable if it fails \\[2pt]
\textbf{This work} & \textbf{yes} & \textbf{yes} & \textbf{yes} &
\textbf{Only lawful possession of the robot} \\
\bottomrule
\end{tabular}
\end{table}


Table~\ref{tab:positioning} sets out the gap from the integrator's side, by asking what each
route to a control interface requires the vendor to have provided first. Every neighbouring
approach either starts once the interface is known, or buys openness at a price a single owned
unit cannot pay; what we describe is the step before all of them.
